\section{Introduction}
\label{sec:introduction}

A useful dynamical Hilbert--Pólya candidate must do more than put its poles on
a circle. Its primitive-orbit data must retain arithmetic information, its
determinant must have the right global divisor, and a natural operator must
survive the passage from finite computations to one fixed infinite object.
Finite graph covers appear to satisfy much of this checklist: nonbacktracking
cycles are exact, voltage holonomy records word order, and Artin--Ihara
theory resolves the zeta function into finite-dimensional representation
sectors \citep{stark1996finite,stark2000partii}.

The most obvious unwanted component is the trivial representation. It carries
the base graph and its branching poles. One might therefore divide the cover
zeta by the base zeta, retain only the nontrivial sectors, and then pass to a
congruence tower. This paper tests that construction before attempting any
spectral fit. The test reveals a three-way obstruction. Canonical aggregation
loses conjugacy-level holonomy; finite representation resolution restores
that information but has the wrong zero density; and the natural amenable
Heisenberg tower develops exact-conductor nontrivial spectrum arbitrarily
close to the trivial spectral edge.

The first obstruction is local and exact. For a primitive base cycle with
holonomy \(g\in G\), put \(o=\ord(g)\) and let \(x\) denote its scalar orbit
weight. The regular cover divided by the trivial Artin factor contributes

\[
(1-x)(1-x^o)^{-|G|/o}.
\]

This standard consequence of the regular representation depends on \(g\)
only through \(o\). To show that the loss affects genuine chronological data,
we give two cyclically reduced primitive words in \(H(\F_7)\) with the same
complete cyclic directed-bigram counts. They are neither cyclic nor
time-reversal copies, but their ordered products are the distinct central
elements \(z^3\) and \(z^2\). The canonical aggregate identifies them because
both elements have order seven. A fixed Schrödinger central character
separates the two phases.

The second obstruction is global. A fixed finite-state, finite-memory graph
model with finitely many positive roofs has a determinant

\[
D(s)=\det\!\left(I-\sum_{j=1}^{r}e^{-s\tau_j}B_j\right).
\]

The determinant is a finite exponential sum. Its zeros in any fixed vertical
strip have count \(O(T)\), by a one-page Jensen argument. The nontrivial
Riemann divisor has count \(\Theta(T\log T)\). Irrationally related finite
roofs remove the simple vertical periodicity of the unit-roof graph zeta, but
they do not change this density mismatch.

The fixed-dimension theorem deliberately leaves growing towers open. We rule
out the uniform Ramanujan/new-sector-gap rescue for the registered natural
Heisenberg congruence tower. The
four-regular Cayley graph of \(H(\Z/3^m\Z)\) contains a primitive
\(3^m\)-dimensional Schrödinger block

\[
H_m=U_m+U_m^*+V_m+V_m^*.
\]

A finite-window Rayleigh vector proves
\(\lambda_{\max}(H_m)\to4\). The central frequency is one, so this block does
not descend from level \(m-1\). It remains after every one-dimensional sector
is deleted. Since the Ramanujan threshold is \(2\sqrt3\), exact-conductor
nontrivial Bass poles eventually leave the graph critical circle and return
toward the branching locations.

The paper makes four scoped contributions.

\begin{enumerate}[leftmargin=2em]
  \item We derive the regular-minus-trivial order formula and state precisely
  which nonabelian information it discards.
  \item We exhibit and independently certify a non-dihedral same-bigram
  chronology witness in \(H(\F_7)\).
  \item We prove a finite-memory zero-density exclusion that covers both
  commensurable and finitely incommensurable roofs.
  \item We prove branch-pole return in a primitive nonabelian sector of the
  Heisenberg congruence tower, so deleting all abelian sectors does not rescue
  this frozen gap mechanism.
\end{enumerate}

The standard graph-covering, finite Heisenberg/Harper, and exponential-sum
inputs are explicitly credited. In particular, the Schrödinger block and its
near-zero-flux edge behavior are prior objects
\citep{dedeo2003spectra,beguin1997harper}; the contribution here is the
conductor-resolved Rayleigh certificate, its frozen Hilbert--Pólya
new-sector-gap application, and the chronology/finite-memory obstruction
synthesis. The result does not cover nuclear infinite-rank
transfer operators, infinitely many roof values, nonlinear spectral changes,
or nonamenable property-\(\tau\) towers. Those exclusions identify the
structural changes that a subsequent construction must make.
